% Standalone problem document, generated from the adjacent README.md.
% Compile with XeLaTeX. Mathematical content is maintained in Markdown.
\documentclass[11pt,a4paper]{article}
\usepackage[fontsize=10.5pt]{scrextend}
\usepackage[margin=22mm,headheight=15pt,headsep=9mm,footskip=11mm]{geometry}
\usepackage{fontspec}
\setmainfont{texgyrepagella}[Extension=.otf,UprightFont=*-regular,BoldFont=*-bold,ItalicFont=*-italic,BoldItalicFont=*-bolditalic]
\setsansfont{texgyreheros}[Extension=.otf,UprightFont=*-regular,BoldFont=*-bold,ItalicFont=*-italic,BoldItalicFont=*-bolditalic]
\setmonofont{texgyrecursor}[Extension=.otf,UprightFont=*-regular,BoldFont=*-bold,ItalicFont=*-italic,BoldItalicFont=*-bolditalic,Scale=MatchLowercase]
\usepackage{amsmath,amssymb}
\usepackage{unicode-math}
\setmathfont{texgyrepagella-math.otf}
\usepackage{xcolor,microtype,enumitem,needspace,fancyhdr}
\usepackage{bookmark}
\definecolor{ink}{HTML}{17344B}
\definecolor{accent}{HTML}{197D80}
\definecolor{muted}{HTML}{596773}
\hypersetup{colorlinks=true,linkcolor=accent,urlcolor=accent,pdftitle={AC-03 - Border
rank of the 3 × 3 matrix
product},pdfauthor={Open Problems in Numerical Linear Algebra}}
\urlstyle{same}
\setlength{\parindent}{0pt}
\setlength{\parskip}{5pt plus 1pt minus 1pt}
\linespread{1.04}
\setlength{\emergencystretch}{3em}
\widowpenalty=10000
\clubpenalty=10000
\displaywidowpenalty=10000
\allowdisplaybreaks[1]
\setlist{nosep,leftmargin=1.5em}
\providecommand{\tightlist}{\setlength{\itemsep}{0pt}\setlength{\parskip}{0pt}}
\providecommand{\pandocbounded}[1]{#1}
\pagestyle{fancy}
\fancyhf{}
\fancyhead[L]{\sffamily\footnotesize\color{muted}OPEN PROBLEMS IN NUMERICAL LINEAR ALGEBRA}
\fancyhead[R]{\sffamily\bfseries\footnotesize\color{accent}AC-03}
\fancyfoot[L]{\parbox[b]{0.9\textwidth}{\sffamily\fontsize{8}{10}\selectfont\color{muted}Editorial ratings; a bounded search cannot certify that a problem remains unsolved.\\Literature check: 2026-09-08}}
\fancyfoot[R]{\sffamily\footnotesize\color{muted}\thepage}
\renewcommand{\headrulewidth}{0.3pt}
\renewcommand{\headrule}{\hbox to\headwidth{\color{accent}\leaders\hrule height \headrulewidth\hfill}}
\makeatletter
\renewcommand{\section}{\Needspace{5\baselineskip}\@startsection{section}{1}{0pt}{12pt plus 2pt minus 2pt}{4pt}{\sffamily\bfseries\large\color{ink}}}
\renewcommand{\subsection}{\Needspace{4\baselineskip}\@startsection{subsection}{2}{0pt}{9pt plus 1pt minus 1pt}{3pt}{\sffamily\bfseries\normalsize\color{ink}}}
\makeatother
\setcounter{secnumdepth}{0}
\begin{document}
{\sffamily\small\color{muted}Arithmetic and complexity\par}
\vspace{3pt}
{\raggedright\hyphenpenalty=10000\exhyphenpenalty=10000\sffamily\bfseries\fontsize{21}{24}\selectfont\color{ink}Border
rank of the \(3\times3\) matrix product\par}
\vspace{7pt}
{\sffamily\small\textbf{Difficulty:} extreme\quad\textbf{Importance:} interesting
to the community\par}
\vspace{4pt}
{\color{accent}\hrule height 0.8pt}
\vspace{6pt}
\textbf{Topic:} approximate bilinear algorithms

\textbf{Status:} open; admitted after a bounded literature search found
no resolution.

\subsection{Problem statement}\label{problem-statement}

In \(\mathbb C^9\otimes\mathbb C^9\otimes\mathbb C^9\), let

\[M_3=\sum_{i,j,k=1}^{3}e_{ij}\otimes e_{jk}\otimes e_{ki}.\]

The tensor rank \(R(T)\) is the minimum number of pure tensors in an
exact sum for \(T\). The border rank \(\underline R(T)\) is the least
\(r\) such that \(T\) is a Euclidean limit of tensors of rank at most
\(r\). Determine \(\underline R(M_3)\). The use of limits makes this a
different invariant from
\href{https://github.com/ajt60gaibb/OpenProblemsInNLA/blob/main/arithmetic-and-complexity/AC-02/README.md}{AC-02}.

\subsection{Why it matters}\label{why-it-matters}

Degenerating bilinear algorithms underlie asymptotic improvements in
matrix multiplication.

\subsection{References and status}\label{references-and-status}

A. Conner, A. Harper, J. M. Landsberg,
\href{https://arxiv.org/abs/1911.07981}{\emph{New lower bounds for
matrix multiplication and det\textsubscript{3}}}, Theorem 1.1, proves
\(\underline R(M_3)\ge17\). J. Alman and B. Li,
\href{https://arxiv.org/abs/2605.21738}{\emph{Asymptotic Rank Speedup
Theorems, Revisited}} (2026), §1, explicitly identifies this exact
border rank as open. Searches for ``3x3 border rank matrix
multiplication 2026'' found no exact determination. \textbf{Admitted: no
resolution located.}
\end{document}
